\section{DATA as a distinguished case}
DATA 3.0 is a rich instance of the general framework. Its complete state includes object theory, target, metatheory, typed verified bank, multiple search frontiers, untrusted proposals, history, learned parameters, resource allocation, candidate/accepted certificates, settlement status, and a monotone round counter \cite{data3}. Its active subsystems are modeled as transformations whose compositions generate a transformation semigroup.

The general theory above interprets DATA in four simultaneous ways:
\begin{enumerate}
\item the instantaneous DATA process is a controlled transition system or IFS-like semigroup action;
\item its cumulative verified history is a SIC;
\item the $P$ route specializes to an ATP, while the combined $P/R/I$ routes form an ALD and the nested dependence/independence route forms an AMLD;
\item the resulting theorem/settlement graph is a rationally weighted maze once exact rational transaction costs are declared.
\end{enumerate}

The DATA-specific intelligence layer can therefore be understood as a collection of maze-navigation policies. Professor scores frontiers; Counselor challenges representations; Picard supplies feedback; Creativity and Dreamer generate candidate moves; Horizon certifies bounded regions; the verified bank adds reusable macro structure. None of these changes the verifier's semantics.

\section{A search funnel for computational accessibility}
The synthesis of the Horizon and Depths programs suggests the following funnel:
\[
\begin{array}{c}
\text{raw staged search dynamics}\\[2pt]
\downarrow\\[2pt]
\text{cumulative verified-history SIC}\\[2pt]
\downarrow\\[2pt]
\text{novelty / quotient / DAG reduction}\\[2pt]
\downarrow\\[2pt]
\text{conservative compilation and shared verified landmarks}\\[2pt]
\downarrow\\[2pt]
\text{verified incumbent bound }B\\[2pt]
\downarrow\\[2pt]
\text{rational branch-and-bound / Dijkstra / A* / exact graph methods}\\[2pt]
\downarrow\\[2pt]
\text{learned ranking, portfolios, reconnaissance, and resurvey}\\[2pt]
\downarrow\\[2pt]
\text{discovery-sensitive control}\\[2pt]
\downarrow\\[2pt]
\text{certificate-bearing fixed point.}
\end{array}
\]
The strongest claim is not that the fixed point becomes easy. It is that each justified transformation can remove a class of states or routes that do not need to be searched for the declared objective, while the discovery layer can learn which surviving transformation family is worth trying next.

\section{Research conjectures and implementation questions}
\begin{conjecture}[Verified-history compression]
On broad formal-mathematics families, searching a canonicalized graph of cumulative verified histories requires asymptotically or practically fewer expansions than searching raw transient histories under matched correctness constraints.
\end{conjecture}

\begin{conjecture}[Compiled-landmark acceleration]
A verified shared bank plus conservative macro compilation reduces rational weighted hitting time on families with repeated substructure, after retrieval and verification overhead are charged exactly.
\end{conjecture}

\begin{conjecture}[Admissible learned lower bounds]
For some nontrivial theorem families, a learned function can be calibrated into a valid rational lower bound $h(v)$ often enough to improve branch-and-bound or A* search while preserving exact minimum-cost guarantees.
\end{conjecture}

\begin{conjecture}[Reflective search improvement]
A meta-ATP can prove bounded or conditional facts about a frozen source ATP/ALD/AMLD that, when deposited into the verified meta-bank, measurably reduce later settlement cost without altering the verifier or certificate semantics.
\end{conjecture}

\begin{conjecture}[Discovery transfer]
There exist formal-mathematics families and verifier-respecting state summaries for which intervention outcomes on earlier solved rungs predict, above matched baselines, which class of intervention increases the probability of a verified novel useful transition on a later hidden rung. The intended claim is empirical and forward-chaining: the policy is frozen before the later rung is exposed.
\end{conjecture}

Key implementation questions include whether cumulative-history canonicalization is cheap enough to pay for itself; how to define rational edge costs that remain meaningful across machines; how much of proof-state equivalence can be checked exactly; whether macro compilation should affect search cost, native proof cost, or both; which meta-theorems are strong enough to improve scheduling rather than merely restating architecture; and whether a system can learn, from earlier settled rungs, when a qualitatively different mathematical mechanism is becoming more valuable than continued local search.
